% Shared methodology: supervised attack pipeline + unsupervised clustering cross-check.
\subsection{Attack methodology}
\label{sec:eval-method}

We evaluate each privacy property with a primary supervised attack and a complementary
unsupervised cross-check. The \emph{supervised} pipeline asks whether a trained
classifier can recover a hidden label from the observables; this is the main privacy
verdict because it gives the adversary labels and therefore more power than a purely
exploratory analysis. The \emph{unsupervised} pipeline asks whether the observables
cluster along the hidden label without any supervision, providing a label-free sanity
check on the same executions.

\subsubsection{Supervised attacks}
For every task we train linear and non-linear classifiers and report the score used by
the unified verdict for that property. Each task is evaluated over \textbf{30 random
seeds} using a \textbf{policy-disjoint train/test split}: the policies used to train the
attacker never appear in the test set, so a high score cannot come from memorizing
individual policies. For each seed we run a label-permutation test and record whether
the AUC is significant at $\alpha{=}0.05$.

The honest headline metrics are the AUC, its distance from chance $|\mathrm{AUC}-0.5|$,
and the permutation $p$-value. We \emph{report} the classifier advantage over random but
do \emph{not} gate on it: on imbalanced test sets a majority-class predictor scores a
large advantage with $F_1\approx 0$, so advantage alone is unreliable. The unified
verdict (\texttt{experiments/core/multiseed\_criterion.js}) is:
%
\begin{quote}
  \textbf{PASS} iff $\mathrm{AUC}_{\text{mean}}\le 0.55$ \textbf{and} the
  significant-seed fraction is $<0.5$ (over 30 policy-disjoint seeds).
\end{quote}
%
Both legs are required: high classification power that is not statistically reproducible
across seeds (or reproducible significance with chance-level AUC) does not constitute a
leak. Subclasses whose median test size is $n<20$ are \emph{underpowered}---reported but
excluded from gating, so a tiny subclass's noise cannot decide the verdict. The result
tables report the AUC selected by the committed verdict file for each task. SPP uses the
calibrated random-forest score for its AUC leg because logistic regression
over-separates on the near-constant SPP features and yields a degenerate AUC.

\paragraph{Positive control.} To show that a chance-level AUC means \emph{no signal}
rather than a \emph{broken attack}, every supervised experiment also runs a synthetic
control: a label deterministically derived from the gas (``\texttt{gas\,>\,median}'')
pushed through the identical pipeline. The control reaches AUC $\approx 0.99$ and
$F_1\approx 0.95$, confirming the attack can detect a real channel when one exists.

\subsubsection{Unsupervised clustering attacks}
The clustering attacker drops the labels entirely and asks whether the gas features
\emph{naturally} group along the hidden property. We run K-Means (with $K$ set to the
number of classes, or $K{=}2$ for policy pairs) and DBSCAN on the on-chain gas features,
and score the recovered clusters against the true labels with the Adjusted Rand Index
(ARI) and Normalized Mutual Information (NMI). Values near $0$ mean the clusters carry no
information about the hidden label; values near $1$ mean full leakage.

PCH clusters all executions globally (no per-layer splitting), so the large-sample
ARI/NMI are read directly. For the PI pairs we additionally guard against two
small-sample failure modes.

\paragraph{Controlling sample-size confounds.} The naive PI pair construction produced
severely imbalanced pairs (one policy with $1$ execution against another with up to
$30$), where K-Means cannot cut the data and ARI/NMI $\approx 0$ is an artifact rather
than evidence of privacy. We therefore rebuild the PI pairs from policies with
$\ge 30$ executions into balanced $30$-vs-$30$ contrasts before clustering. To rule out
residual finite-sample inflation, we add a \emph{permutation baseline}: we shuffle the
pair labels $200$ times and recompute the clustering NMI, giving the NMI distribution
under zero true signal. The debiased NMI is the observed value minus this baseline, and
the empirical $p$-value is the fraction of permutations whose NMI meets or exceeds the
observed; a pair leaks only if its debiased NMI is positive \emph{and} significant
($p<0.05$).

\paragraph{Dataset.} Both attack families run on the same fixed-circuit, commitment-only,
access-gas-padded \emph{Polygon~Amoy} dataset: the supervised verdicts read
\texttt{experiments/multiseed\_results/polygonAmoy/} (numbers of record), and the
clustering reads the same executions exported to \texttt{python/data/polygonAmoy/}. The
supervised and unsupervised attacks therefore constitute two independent attack families
applied to one dataset, rather than two datasets. Per-task sample sizes vary and are
listed in the result tables (e.g.\ $2005$ DUO executions; $133$ balanced $30$-vs-$30$ PI
pairs); the PI debiasing averages each observed NMI over $20$ K-Means seeds against a
$200$-permutation baseline. Provenance for every table and figure is listed in
\texttt{paper/README.md}.
